\documentclass[12pt,conference,onecolumn]{IEEEtran}

\usepackage[hidelinks]{hyperref}

\title{Spy car}
\author{%
\IEEEauthorblockN{Andrew Kabatsky}\IEEEauthorblockA{Science \& Engineering\\Manalapan High School\\Englishtown, NJ\\\href{mailto:426akabatsky@frhsd.com}{426akabatsky@frhsd.com}}\and
\IEEEauthorblockN{Sharone Krasnopolsky}\IEEEauthorblockA{Science \& Engineering\\Manalapan High School\\Englishtown, NJ\\\href{mailto:426skrasnopolsky@frhsd.com}{426skrasnopolsky@frhsd.com}}\and
\IEEEauthorblockN{Grace Nealon}\IEEEauthorblockA{Science \& Engineering\\Manalapan High School\\Englishtown, NJ\\\href{mailto:426gnealon@frhsd.com}{426gnealon@frhsd.com}}}
\date{January 28, 2026}

\newcommand{\keywords}{Traffic Safety Division, Monmouth County Engineering, spy car, road study, Raspberry Pi, 802.11, wi-fi, video}

\usepackage{hyperref}
\makeatletter
\AtBeginDocument{
\hypersetup{%
pdftitle={\@title},
pdfauthor={Grace Nealon},
pdfkeywords={\keywords}}}
\makeatother

\begin{document}
\maketitle 

\begin{abstract}
For the duration of this semester, we will design and build a remotely controlled car that utilizes a camera to wirelessly relay a live video feed from the car back to the user. The project will consist of three main parts: the car, the controller, and the camera. We plan to experiment with different chassis and use a Logitech F310 to control our rover. We want to create smooth communication between the car and controller through the use of a WiFi connection and a Raspberry Pi setup. We also want to relay live video and transmit data using serial communication.

This project will allow us to study different aspects of robotics, programming, and project design. We will learn to use microcontrollers to command different electronic components, and our goal is to obtain wireless communication between devices. We will combine both software and hardware aspects in this project.
\end{abstract}

\begin{IEEEkeywords}
\keywords
\end{IEEEkeywords}

\end{document}
